\section{Verification purpose}

\textbf{Case ID:} \texttt{\CaseID}. This case verifies PER-MW firebrand generation and empirical Lagrangian transport for surface fire.  The verification question is whether the dominant emitting cell's count and log-distance distribution agree with the independent GUIDE targets.

\section{Mathematical formulation and numerical implementation}

For the selected source cell, firebrands are generated in proportion to fire power. Empirical transport produces positive distances \(d_i\). A lognormal response is summarized by \(\mu=\mathrm{mean}[\ln d_i]\) and \(\sigma=\mathrm{std}[\ln d_i]\). The postprocessor groups the spotting log by source indices before calculating these quantities.

\section{Derivation of expected behavior}

The configured experiment predicts 8450 firebrands from the dominant cell and log-distance parameters \(\mu=2.42\) and \(\sigma=1.31\). Grouping by source is essential: combining cells would mix different emission histories and would not test the intended distribution.

\section{Verification metrics and acceptance criteria}

Dominant-cell count detects generation-rate or duration errors and uses 15 percent relative tolerance to accommodate stochastic sampling. The mean and standard deviation of \(\ln d\) test location and width of the transport distribution and each use absolute tolerance 0.15. All three components require complete spotting logs and must pass.

\section{Simulation configuration and predicted outcome}

The case uses Fuel Model 3, flat terrain, 5 mph wind, PER-MW generation, the EMPIRICAL distance model, LAGRANGIAN accumulation, DIRECT ignition, and deterministic random-seed handling. Spotting statistics output is required.

\section{Actual simulation results}

The required evidence is generated only from case-local ELMFIRE outputs. The current postprocessing status is \textbf{\VerificationStatus}; \PassingMetricCount{} of \MetricCount{} required metric rows pass across \VariantCount{} variant(s). A missing output is reported as NOT RUN or NOT EVALUABLE and is not interpreted as verification evidence.


\begin{figure}[H]
  \centering
  \IfFileExists{../figures/whole_domain_result.pdf}{%
    \includegraphics[width=0.82\textwidth,height=0.58\textheight,keepaspectratio]{whole_domain_result.pdf}%
  }{\fbox{\parbox{0.85\textwidth}{No case-local whole-domain ELMFIRE raster is available.}}}
  \caption{Final whole-domain ELMFIRE time of arrival for the displayed variant, in seconds. The figure annotation records the selected case-local GeoTIFF path and coordinate reference system; the generated metrics table, not visual appearance, determines acceptance.}
\end{figure}


\section{Calculated metrics and verification decision}

\begin{center}
\begin{tabular}{@{}llll@{}}
\toprule
Variant / metric & Expected & Calculated & Status \\
\midrule
\MetricRows
\bottomrule
\end{tabular}
\end{center}

The authoritative overall decision is \textbf{\VerificationStatus} and is generated by \path{scripts/postprocess.py} from \path{outputs/metrics.json}; it is not recomputed in this report.

\section{Technical interpretation}

Count-only failure points toward generation or residence time; distribution-only failure points toward sampling, units, or transport parameterization. No scientific conclusion should be drawn when the status is NOT RUN or NOT EVALUABLE.

\section{Scope and limitations}

The metrics cover the dominant source cell and one forcing condition. They verify implementation against the prescribed statistical target, not physical firebrand observations. This case provides implementation verification evidence only.
